\documentclass{article}

% NeurIPS 2026 workshop style
\usepackage[final,nonatbib]{neurips_2024}  % use neurips_2024 as base; replace with neurips_2026 when available

\usepackage[utf8]{inputenc}
\usepackage[T1]{fontenc}
\usepackage{hyperref}
\usepackage{url}
\usepackage{booktabs}
\usepackage{amsfonts}
\usepackage{amsmath}
\usepackage{amssymb}
\usepackage{nicefrac}
\usepackage{microtype}
\usepackage{graphicx}
\usepackage{xcolor}
\usepackage{multirow}
\usepackage{subcaption}
\usepackage{natbib}

\bibliographystyle{abbrvnat}

\newcommand{\Kcal}{\mathcal{K}}
\newcommand{\nats}{\,\text{nats}}

\title{Phase Boundaries of Circuit Discovery in Neural Networks}

\author{
  Mihir Sahasrabudhe \\
  School of Information Sciences \\
  University of Illinois Urbana-Champaign \\
  \texttt{mihirrs2@illinois.edu}
}

\begin{document}

\maketitle

\begin{abstract}
When transformers learn surjective mappings $(B, z) \to A$ with $K$-fold ambiguity, training exhibits a metastable plateau at loss $\log K$ followed by a sharp phase transition---the \emph{disambiguation lag}.
We map the $(K, \eta)$ phase diagram under AdamW and identify discoverable, nucleation, and unreachable zones.
The phase boundary scales as $\eta^*(K) \propto K^{-0.83}$ ($R^2 = 0.96$), and within the discoverable region the waiting time scales as $\tau \propto K^{1.31}$ ($R^2 = 0.98$).
The power-law scaling $\tau \propto K^\alpha$ holds across three architectures (Transformer $\alpha{=}1.70$, Gated MLP $\alpha{=}1.62$, RNN $\alpha{=}1.47$; all $R^2 {>} 0.96$).
A batch-size control disentangles gradient noise from step size, showing that \emph{increased} noise delays the transition---opposite to Kramers-type thermal escape.
Post-hoc weight analysis reveals that the plateau is a random walk in weight space (direction consistency ${\approx}0$) that snaps into coherent descent at $\tau$; a synergistic $z$-dependence signal precedes the loss transition by ${\sim}50\%$ of $\tau$ across all $K$.
These results identify an empirical optimization-geometric constraint on compositional generalization at the tested model scale.
\end{abstract}


%==============================================================================
\section{Introduction}
%==============================================================================

Large language models fail at tasks requiring contextual selection among competing associations.
The reversal curse~\citep{berglund2024reversal} and factorization curse~\citep{kitouni2024factorization} show that variable-binding operations resist learning even when data and capacity suffice.
Mitigations like bidirectional augmentation~\citep{lv2024bico} and reverse training~\citep{golovneva2024reverse} address the symptom, but the underlying optimization dynamics remain poorly understood.

We propose that the barrier is \emph{optimization-geometric}: the target circuit occupies a narrow basin whose accessibility shrinks with the combinatorial complexity $K$ of the binding.
We construct a synthetic task with a single control parameter $K$ and map the $(K, \eta)$ phase diagram under AdamW at a fixed model scale.
We report five findings: (1)~the plane partitions into discoverable, nucleation, and unreachable zones separated by $\eta^*(K) \propto K^{-0.83}$ (\S\ref{sec:phase-diagram}); (2)~$\tau \propto K^{1.31}$ in the discoverable region (\S\ref{sec:dynamics}); (3)~the power-law holds across three architectures (\S\ref{sec:architecture}); (4)~a batch-size control cleanly disentangles noise and step size, falsifying thermal escape (\S\ref{sec:batch-size}); (5)~a synergistic $z$-dependence signal leads the loss transition by ${\sim}50\%$ of $\tau$, while weight displacement analysis reveals the plateau as a random walk that snaps into coherent descent at the transition (\S\ref{sec:discussion}).


%==============================================================================
\section{The Disambiguation Task}
\label{sec:task}
%==============================================================================

We define a surjective map $f: \mathcal{A} \to \mathcal{B}$ with $|\mathcal{B}| = 1000$ base strings and constant fiber size $|f^{-1}(b)| = K$.
Each fiber is indexed by a selector $z \in \{1, \ldots, K\}$, making $(b, z) \mapsto a$ one-to-one.
The model receives $[\texttt{BOS}, B, \texttt{SEP}, z, \texttt{SEP}]$ and predicts the 4-character target $A$ autoregressively.
The conditional entropy $H(A \mid B) = \log K$; $H(A \mid B, z) = 0$.
Two-layer linear networks converge to $\log K$ permanently with zero $z$-dependence, confirming that disambiguation requires multiplicative capacity.

We train a 4-layer Transformer ($d_\text{model} = 128$, 4 heads, $d_\text{mlp} = 512$; ${\sim}600$K parameters) with AdamW~\citep{loshchilov2019adamw}, batch size 128, cosine warmup, via TransformerLens~\citep{nanda2022transformerlens}.
Training budgets are 10k--50k steps (up to 100k for $K{=}36$).
The key diagnostic is the \emph{$z$-shuffle gap}: permuting $z$ across the batch while holding $(B, A)$ fixed measures $\Delta_z = \mathcal{L}_{z\text{-shuffle}} - \mathcal{L}_\text{clean}$; during the plateau $\Delta_z = 0$, after the transition $\Delta_z \gg \log K$.


%==============================================================================
\section{The Phase Diagram}
\label{sec:phase-diagram}
%==============================================================================

\subsection{Two-phase dynamics}

Every run with $K > 1$ exhibits two sharply separated phases.
In the \emph{marginal phase}, loss drops to $\approx \log K\nats$ within a few hundred steps ($\Delta_z = 0$).
After a prolonged plateau, the \emph{conditional phase} begins: loss undergoes a sharp sigmoidal decline as the model discovers the binding circuit.

\subsection{The $(K, \eta)$ phase diagram}
\label{sec:phase-boundary}

We conduct 51 training runs spanning $K \in \{5, 10, 20, 36\}$, multiple $\eta$ per $K$, and 3 seeds per pair.
Runs are classified as \emph{success} (loss below $\log K$) or \emph{failure} (no transition within the training budget).%
\footnote{The ``unreachable'' designation is budget-relative; a long-enough run might transition. However, the $\eta{=}2{\times}10^{-3}$, $K{=}20$ run had not fully converged after 50k steps ($\tau \approx 22{,}500$), suggesting the boundary is meaningful.}

\begin{figure}[t]
  \centering
  \includegraphics[width=\textwidth]{outputs/paper_figures/fig_phase_diagram.pdf}
  \caption{\textbf{The $(K, \eta)$ phase diagram.} Three zones: \textcolor{green!60!black}{\textbf{discoverable}} (all seeds converge), \textcolor{orange!80!black}{\textbf{nucleation}} (partial convergence), \textcolor{red!70!black}{\textbf{unreachable}} (no convergence within budget). Dashed: $\eta^*(K) = 0.048 \cdot K^{-0.83}$.}
  \label{fig:phase-diagram}
\end{figure}

Figure~\ref{fig:phase-diagram} reveals three zones.
The boundary is well-described by $\eta^*(K) = 0.048 \cdot K^{-0.83}$ ($R^2 = 0.96$, $\delta = 0.83 \pm 0.16$), where $\eta^*$ is the geometric mean of the highest all-succeed and lowest any-fail learning rates (Appendix~\ref{app:matrix}).

\subsection{Dynamics inside the discoverable region}
\label{sec:dynamics}

At fixed $\eta = 10^{-3}$, we measure $\tau$ across $K \in \{3, 5, 7, 10, 13, 17, 20, 25, 30, 36\}$:
\begin{equation}
  \tau \propto K^{1.31}, \quad R^2 = 0.983, \quad \text{LOO range: } [1.27, 1.36]
  \label{eq:tau}
\end{equation}
Near the boundary, $\tau$ diverges: at $K{=}36$, increasing $\eta$ from $10^{-3}$ to $2{\times}10^{-3}$ increases $\bar{\tau}$ by $6.3\times$.
The transition is stochastic near the boundary (seed-dependent, with initiation without completion), consistent with a nucleation picture.


\begin{figure}[t]
  \centering
  \begin{subfigure}[t]{0.48\textwidth}
    \centering
    \includegraphics[width=\textwidth]{outputs/paper_figures/fig_phase_boundary_trajectories.pdf}
    \caption{Loss trajectories for $K{=}36$ near the boundary.}
    \label{fig:trajectories}
  \end{subfigure}
  \hfill
  \begin{subfigure}[t]{0.48\textwidth}
    \centering
    \includegraphics[width=\textwidth]{outputs/paper_figures/fig_landauer_dense.pdf}
    \caption{$\tau$ vs.\ $K$ at $\eta = 10^{-3}$ with power-law fit $\tau \propto K^{1.31}$.}
    \label{fig:tau-scaling}
  \end{subfigure}
  \caption{\textbf{Dynamics near and within the discoverable region.}}
  \label{fig:dynamics}
\end{figure}


%==============================================================================
\section{Controls and Falsifications}
\label{sec:controls}
%==============================================================================

\paragraph{Not gradient cancellation.}
Varying the number of interfering groups $G$ from 1 to 1000 at $K{=}20$ yields exponent $\beta \approx 0.02$ (predicted ${\approx}1$): $\tau$ is invariant to sharing structure.

\paragraph{Not Kramers-type thermal escape.}
If the plateau were a free-energy trap, larger $\eta$ (more noise) should facilitate escape.
The opposite holds: $\tau$ increases monotonically with $\eta$ (at $K{=}20$: $\tau = 2{,}100$ at $\eta = 3{\times}10^{-4}$; $22{,}500$ at $2{\times}10^{-3}$; diverges at $5{\times}10^{-3}$; Appendix~\ref{app:lrsweep}).

\paragraph{Batch-size disentanglement.}
\label{sec:batch-size}
A potential objection is that $\eta$ conflates noise scale and step size~\citep{smith2018dont}.
We sweep $B \in \{64, 128, 256, 512\}$ at fixed $\eta = 10^{-3}$, varying noise ($\propto 1/\sqrt{B}$) without changing step size.
At $K{=}20$: $B{=}64$ fails; $B{=}128 \to \tau{=}6{,}500$; $B{=}256 \to 3{,}600$; $B{=}512 \to 1{,}850$.
At $K{=}36$: $B{=}64$ fails; $B{=}128 \to 22{,}900$; $B{=}512 \to 5{,}300$ (Appendix~\ref{app:batchsize}).
Less noise $\Rightarrow$ faster transition---opposite to Kramers.
At matched SGD temperature $T{=}\eta/B$, larger step size converges $2.2\times$ faster ($B{=}256, \eta{=}10^{-3}$: $\tau{=}3{,}600$ vs.\ $B{=}128, \eta{=}5{\times}10^{-4}$: $\tau{=}7{,}750$).

\paragraph{Architecture universality.}
\label{sec:architecture}
All three nonlinear architectures (Transformer, LSTM, Gated MLP) exhibit two-phase dynamics (Figure~\ref{fig:architectures}a).
Critically, all display power-law scaling $\tau \propto K^\alpha$: Transformer $\alpha{=}1.70$ ($R^2{=}0.97$, $K {\in} [3, 30]$), Gated MLP $\alpha{=}1.62$ ($R^2{=}0.97$, $K {\in} [10, 50]$), RNN $\alpha{=}1.47$ ($R^2{=}0.98$, $K {\in} [10, 100]$).
The universality of the power-law form---despite different exponents and prefactors---suggests the phase boundary reflects task geometry rather than architecture-specific properties.

\begin{figure}[t]
  \centering
  \begin{subfigure}[t]{0.48\textwidth}
    \centering
    \includegraphics[width=\textwidth]{outputs/paper_figures/fig_arch_comparison.pdf}
    \caption{Loss trajectories at $K{=}20$ for three architectures.}
    \label{fig:arch-comp}
  \end{subfigure}
  \hfill
  \begin{subfigure}[t]{0.48\textwidth}
    \centering
    \includegraphics[width=\textwidth]{outputs/paper_figures/fig_cross_arch_tau.pdf}
    \caption{$\tau$ vs.\ $K$ on log-log axes with power-law fits.}
    \label{fig:cross-arch-tau}
  \end{subfigure}
  \caption{\textbf{Architecture controls.} (a)~Two-phase dynamics are universal. (b)~$\tau \propto K^\alpha$ holds for each architecture. Linear models converge to $\log K$ permanently (Appendix~\ref{app:linear}).}
  \label{fig:architectures}
\end{figure}

\begin{figure}[t]
  \centering
  \begin{subfigure}[t]{0.48\textwidth}
    \centering
    \includegraphics[width=\textwidth]{outputs/paper_figures/fig_synergy_proxy_panel_b.pdf}
    \caption{Synergy onset precedes loss onset across all $K$.}
    \label{fig:synergy-leads}
  \end{subfigure}
  \hfill
  \begin{subfigure}[t]{0.48\textwidth}
    \centering
    \includegraphics[width=\textwidth]{outputs/paper_figures/fig_displacement_panel_d.pdf}
    \caption{Direction consistency of weight updates.}
    \label{fig:direction-consistency}
  \end{subfigure}
  \caption{\textbf{Post-hoc analyses.} (a)~$\Delta_z$ onset vs.\ loss onset; all points below identity ($\Delta t_\text{lead} {\propto} K^{2.13}$, $R^2{=}0.99$). (b)~Cosine similarity of consecutive weight deltas: the plateau is a random walk ($\cos\theta {\approx} 0$) that snaps into directed descent at $\tau$.}
  \label{fig:posthoc}
\end{figure}

\paragraph{Synergy cascade and circuit nucleation.}
The $z$-shuffle gap $\Delta_z$ rises \emph{before} the loss transition across all $K$ values (Appendix~\ref{app:synergy}).
The lead time scales as $\Delta t_\text{lead} \propto K^{2.13}$ ($R^2 {=} 0.99$), with a mean lead fraction of $0.52 \pm 0.14$ of $\tau_\text{loss}$ (Figure~\ref{fig:posthoc}a).
Per-head probing at $K{=}20$ reveals a three-stage cascade: head L1H3 begins attending to $z$ at step 500, the aggregate $\Delta_z$ exceeds $0.1\nats$ at step 2{,}750, and loss drops at step 8{,}700.
The individual circuit signal precedes the aggregate information-theoretic signal by ${\sim}2{,}200$ steps.


%==============================================================================
\vspace{-3pt}
\section{Discussion}
\label{sec:discussion}
\vspace{-2pt}
%==============================================================================

\paragraph{Connection to the reversal curse.}
In the reversal curse, ``$B \to A$'' has fan-in $K$ while ``$A \to B$'' is one-to-one; our results predict the backward mapping requires ${\sim}K^{1.3}$--$1.7\times$ more training.
The compositionality gap~\citep{press2023measuring} and autoregressive biases~\citep{mccoy2023embers} may interact with the phase boundary in natural-language settings.

\paragraph{Hidden progress and weight dynamics.}
Both grokking~\citep{power2022grokking, nanda2023progress} and our task exhibit metastable plateaus before sharp transitions, but our task provides a continuous control parameter ($K$) and quantitative scaling laws.
Weight displacement analysis (Appendix~\ref{app:displacement}) confirms the hidden progress picture of~\citet{barak2022hidden}: during the plateau, $\|W(t{+}\Delta) - W(t)\|_F$ is large and approximately constant (${\sim}11$ per 200-step interval at $K{=}20$), but consecutive displacement directions are \emph{uncorrelated} ($\cos\theta \approx 0.04$).
At $\tau$, direction consistency snaps to ${\sim}0.8$ (Figure~\ref{fig:posthoc}b): the transition is not the onset of weight movement but the onset of \emph{coherent} movement---a symmetry-breaking event where the optimizer finds a descent channel after diffusing on a nearly-flat manifold~\citep{chen2024hidden}.

\paragraph{Connections to recent theory.}
We tested predictions from six theoretical frameworks against our checkpoint and training-history data (Appendices~\ref{app:synergy}--\ref{app:adamw}).
The synergy-as-order-parameter framework~\citep{aguilera2024synergy} is strongly supported: $\Delta_z$ onset systematically precedes loss onset with a clean power-law lead time (\S\ref{sec:controls}).
Weight norm dynamics partially confirm deep grokking predictions~\citep{gromov2024grokking}: while total $\|W\|_F$ grows monotonically, the layer containing the nucleating head shows non-monotonic behavior (Appendix~\ref{app:weight-norms}).
OV circuit spectra~\citep{tigges2024cdt} reveal distributed effective rank changes; L1H3 maintains high effective rank throughout, operating as a low-norm precision instrument rather than a rank-1 projector (Appendix~\ref{app:spectrum}).
The AdamW composite timescale $\tau_\text{adamw} {=} B/(\eta\lambda D)$~\citep{lyu2025power} anti-correlates with $\tau_\text{trans}$, and within-group representation variance \emph{increases} at the transition---opposite to neural collapse predictions~\citep{rangamani2025neural}---because the model must differentiate $z$ within each $B$-group (Appendix~\ref{app:neural-collapse}).
Cumulative gradient waste during the plateau scales as $\tau^{2.14}$ ($R^2 {=} 0.99$; Appendix~\ref{app:gradient}).

\paragraph{Limitations.}
The boundary fit uses four $K$ values ($<$1 decade; $\delta = 0.83 \pm 0.16$); K-sweep spans $K {\in} [3, 36]$.
All experiments use AdamW at one model scale (${\sim}600$K params) on a synthetic task; other optimizers, scales, and natural-language data may shift the boundary.

\bibliography{paper_workshop_v1}

\appendix
\section{Full Phase Boundary Outcome Matrix}
\label{app:matrix}

\begin{table}[h]
\centering
\caption{\textbf{Full outcome matrix.} Success rate (of 3 seeds) and mean $\tau$ per $(K, \eta)$. ``---'' = no convergence.}
\label{tab:full-matrix}
\vspace{4pt}
\small
\begin{tabular}{@{}cccc@{}}
\toprule
$K$ & $\eta$ & Rate & Mean $\tau$ \\
\midrule
5  & $5.0 \times 10^{-3}$ & 3/3 & 2{,}850 \\
5  & $7.0 \times 10^{-3}$ & 3/3 & 3{,}250 \\
5  & $1.0 \times 10^{-2}$ & 3/3 & 2{,}183 \\
5  & $1.5 \times 10^{-2}$ & 0/3 & --- \\
5  & $2.0 \times 10^{-2}$ & 0/3 & --- \\
\midrule
10 & $3.0 \times 10^{-3}$ & 3/3 & 5{,}950 \\
10 & $5.0 \times 10^{-3}$ & 3/3 & 9{,}867 \\
10 & $7.0 \times 10^{-3}$ & 3/3 & 8{,}533 \\
10 & $1.0 \times 10^{-2}$ & 1/3 & 5{,}300 \\
\midrule
20 & $2.0 \times 10^{-3}$ & 3/3 & 5{,}750 \\
20 & $3.0 \times 10^{-3}$ & 3/3 & 30{,}300 \\
20 & $5.0 \times 10^{-3}$ & 2/3 & 22{,}425 \\
20 & $7.0 \times 10^{-3}$ & 1/3 & 24{,}850 \\
\midrule
36 & $1.0 \times 10^{-3}$ & 3/3 & 10{,}867 \\
36 & $2.0 \times 10^{-3}$ & 2/3 & 68{,}350 \\
36 & $3.0 \times 10^{-3}$ & 2/3 & 29{,}225 \\
36 & $5.0 \times 10^{-3}$ & 0/3 & --- \\
\bottomrule
\end{tabular}
\end{table}

\section{Batch-Size Sensitivity}
\label{app:batchsize}

\begin{table}[h]
\centering
\caption{\textbf{Batch-size sensitivity} at fixed $\eta = 10^{-3}$. ``---'' = no convergence.}
\label{tab:batch-size}
\vspace{4pt}
\small
\begin{tabular}{@{}cccccc@{}}
\toprule
$K$ & $B$ & Noise $\propto 1/\sqrt{B}$ & $T{=}\eta/B$ & $\tau$ (steps) & Conv? \\
\midrule
20 & 64  & 0.125 & $1.56 \times 10^{-5}$ & ---    & No \\
20 & 128 & 0.088 & $7.81 \times 10^{-6}$ & 6{,}500  & Yes \\
20 & 256 & 0.063 & $3.91 \times 10^{-6}$ & 3{,}600  & Yes \\
20 & 512 & 0.044 & $1.95 \times 10^{-6}$ & 1{,}850  & Yes \\
\midrule
36 & 64  & 0.125 & $1.56 \times 10^{-5}$ & ---    & No \\
36 & 128 & 0.088 & $7.81 \times 10^{-6}$ & 22{,}900 & Yes \\
36 & 256 & 0.063 & $3.91 \times 10^{-6}$ & 8{,}650  & Yes \\
36 & 512 & 0.044 & $1.95 \times 10^{-6}$ & 5{,}300  & Yes \\
\bottomrule
\end{tabular}
\end{table}

\section{Dense K-Sweep Results}
\label{app:ksweep}

\begin{table}[h]
\centering
\caption{\textbf{Dense K-sweep} at $\eta = 10^{-3}$, seed 42.}
\label{tab:ksweep}
\vspace{4pt}
\small
\begin{tabular}{@{}cccc@{}}
\toprule
$K$ & $\log K$ (nats) & $\tau$ (steps) & Transition window \\
\midrule
3  & 1.10 & 200  & 300--500 \\
5  & 1.61 & 400  & 700--1{,}100 \\
7  & 1.95 & 650  & 800--1{,}450 \\
10 & 2.30 & 850  & 1{,}600--2{,}450 \\
13 & 2.56 & 1{,}350 & 1{,}800--3{,}150 \\
17 & 2.83 & 2{,}700 & 2{,}700--5{,}400 \\
20 & 3.00 & 3{,}400 & 3{,}550--6{,}950 \\
25 & 3.22 & 4{,}100 & 4{,}700--8{,}800 \\
30 & 3.40 & 4{,}700 & 5{,}700--10{,}400 \\
36 & 3.58 & 6{,}350 & 7{,}300--13{,}650 \\
\bottomrule
\end{tabular}
\end{table}

\section{Learning Rate Sweep at $K{=}20$}
\label{app:lrsweep}

\begin{table}[h]
\centering
\caption{\textbf{LR sweep at $K{=}20$.} $\tau$ increases monotonically with $\eta$; $\eta = 5 \times 10^{-3}$ does not transition within 50k steps.}
\label{tab:lrsweep}
\vspace{4pt}
\small
\begin{tabular}{@{}ccccc@{}}
\toprule
$\eta$ & $\tau$ (steps) & $\eta \cdot \tau$ & $Q_\text{transition}$ & Conv? \\
\midrule
$3 \times 10^{-4}$ & 2{,}100 & 0.63 & 11.5 & Yes \\
$5 \times 10^{-4}$ & 2{,}200 & 1.10 & 13.9 & Yes \\
$1 \times 10^{-3}$ & 3{,}400 & 3.40 & 22.5 & Yes \\
$2 \times 10^{-3}$ & 7{,}550 & 15.10 & 44.5 & Yes \\
$5 \times 10^{-3}$ & --- & --- & --- & No \\
\bottomrule
\end{tabular}
\end{table}

\section{Linear Baseline}
\label{app:linear}

\begin{figure}[h]
  \centering
  \includegraphics[width=0.5\textwidth]{outputs/paper_figures/fig_linear.pdf}
  \caption{Two-layer linear network at $K \in \{10, 20, 36\}$. Loss converges to $\log K$ permanently; $z$-gap is zero throughout.}
  \label{fig:linear}
\end{figure}


%==============================================================================
% Post-hoc experiments connecting to recent theoretical frameworks
%==============================================================================

\section{Synergy Cascade Analysis}
\label{app:synergy}

The $z$-shuffle gap $\Delta_z = \mathcal{L}_{z\text{-shuffle}} - \mathcal{L}_\text{clean}$ serves as a proxy for synergistic information~\citep{aguilera2024synergy}.
We measure onset times for both $\Delta_z > 0.1\nats$ and loss $< 0.9 \log K$ across $K \in \{3, \ldots, 36\}$.

\begin{figure}[h]
  \centering
  \includegraphics[width=\textwidth]{outputs/paper_figures/fig_synergy_proxy.pdf}
  \caption{\textbf{Synergy leads loss.} (a)~Normalized loss (solid) and $\Delta_z$ (dashed) for selected $K$. (b)~Synergy onset vs.\ loss onset; all points below identity. (c)~Lead time vs.\ $K$ on log-log axes: $\Delta t_\text{lead} \propto K^{2.13}$ ($R^2 = 0.99$).}
  \label{fig:synergy-full}
\end{figure}

\begin{table}[h]
\centering
\caption{\textbf{Synergy vs.\ loss onset.} Synergy onset precedes loss onset for all $K$ with mean lead fraction $0.52 \pm 0.14$.}
\label{tab:synergy}
\vspace{4pt}
\small
\begin{tabular}{@{}cccccc@{}}
\toprule
$K$ & $\log K$ & $\tau_\text{loss}$ & $\tau_\text{synergy}$ & Lead (steps) & Lead$/\tau$ \\
\midrule
3  & 1.10 & 500    & 400    & 100    & 0.20 \\
5  & 1.61 & 1{,}100  & 700    & 400    & 0.36 \\
7  & 1.95 & 1{,}300  & 650    & 650    & 0.50 \\
10 & 2.30 & 2{,}350  & 1{,}250  & 1{,}100  & 0.47 \\
13 & 2.56 & 3{,}000  & 1{,}300  & 1{,}700  & 0.57 \\
17 & 2.83 & 5{,}900  & 2{,}700  & 3{,}200  & 0.54 \\
20 & 3.00 & 8{,}700  & 3{,}800  & 4{,}900  & 0.56 \\
25 & 3.22 & 14{,}550 & 5{,}550  & 9{,}000  & 0.62 \\
30 & 3.40 & 24{,}500 & 7{,}800  & 16{,}700 & 0.68 \\
36 & 3.58 & 36{,}750 & 11{,}600 & 25{,}150 & 0.68 \\
\bottomrule
\end{tabular}
\end{table}

Per-head analysis at $K{=}20$ reveals L1H3 begins attending to $z$ at step 500---2{,}250 steps before the aggregate $\Delta_z$ onset (step 2{,}750) and 8{,}200 steps before the loss transition.

\section{Hidden Progress in Weight Space}
\label{app:displacement}

We compute weight displacement $\|W(t{+}\Delta) - W(t)\|_F$ between consecutive checkpoints (every 200 steps) for $K \in \{10, 20, 36\}$ and the cosine similarity of consecutive displacement vectors to measure direction consistency~\citep{barak2022hidden, chen2024hidden}.

\begin{figure}[h]
  \centering
  \includegraphics[width=\textwidth]{outputs/paper_figures/fig_weight_displacement.pdf}
  \caption{\textbf{Weight displacement dynamics.} (a)~Total displacement per step (log scale). (b)~Relative displacement (normalized by $\|W(t)\|$). (c)~Per-component displacement at $K{=}20$: attention spikes highest at $\tau$. (d)~Direction consistency: near-zero during plateau (random walk), snaps to ${\sim}0.8$ at transition.}
  \label{fig:displacement-full}
\end{figure}

\begin{table}[h]
\centering
\caption{\textbf{Phase-wise displacement statistics} at $K{=}20$ ($\tau = 10{,}750$).}
\label{tab:displacement}
\vspace{4pt}
\small
\begin{tabular}{@{}lccc@{}}
\toprule
Phase & Displacement & Relative disp. & $\cos\theta$ \\
\midrule
Plateau ($t < \tau/2$) & $11.27 \pm 4.22$ & 0.128 & 0.036 \\
Pre-transition & $9.44 \pm 0.36$ & 0.079 & $-0.179$ \\
Post-transition ($t > \tau$) & $3.93 \pm 4.71$ & 0.030 & 0.397 \\
\bottomrule
\end{tabular}
\end{table}

The negative direction consistency in the pre-transition window ($\cos\theta \approx -0.18$) suggests an oscillatory regime where the optimizer alternates directions before finding the descent channel.

\section{Weight Norm Dynamics}
\label{app:weight-norms}

Total Frobenius norm $\|W\|_F$ across training for $K \in \{3, 5, 10, 20, 36\}$, sampled every 500 steps~\citep{gromov2024grokking}.

\begin{figure}[h]
  \centering
  \includegraphics[width=\textwidth]{outputs/paper_figures/fig_weight_norms.pdf}
  \caption{\textbf{Weight norm dynamics.} (a)~Total norm grows monotonically; higher $K$ produces larger norms. (b)~Per-component norms at $K{=}20$: unembed shows the sharpest rise at $\tau$. (c)~Per-layer norms: Layer~1 (nucleating layer) is the only layer with non-monotonic behavior (peak then decay).}
  \label{fig:weight-norms-full}
\end{figure}

Total weight norms grow monotonically---no global peak-and-reorganize as predicted by the deep grokking framework. However, Layer~1 (containing nucleating head L1H3) peaks at $\tau$ then decays, showing localized non-monotonic dynamics. Peak/init ratio scales from $1.13\times$ (K=3) to $2.53\times$ (K=36).

\section{OV Circuit Spectrum}
\label{app:spectrum}

For each attention head, we compute $W_\text{OV} = W_V W_O$ and track its singular value spectrum across training~\citep{tigges2024cdt}. Effective rank is the participation ratio $(\sum \sigma_i)^2 / \sum \sigma_i^2$.

\begin{figure}[h]
  \centering
  \includegraphics[width=\textwidth]{outputs/paper_figures/fig_circuit_spectrum.pdf}
  \caption{\textbf{OV circuit spectrum} at $K{=}20$. (a)~Effective rank heatmap; broad rank reduction at $\tau$, strongest in Layers~0 and~3. (b)~$\sigma_1/\sigma_2$: L0H0 shows an early transient spike; L1H3 remains flat. (c)~Nuclear norm: L1H3 has persistently lower nuclear norm than all other heads. (d)~L1H3 effective rank across $K$ values.}
  \label{fig:spectrum-full}
\end{figure}

L1H3 maintains high effective rank (${\sim}26$--$28$) and low $\sigma_1/\sigma_2$ ratio (${\sim}1.05$--$1.09$) throughout training, indicating it operates as a distributed, low-norm circuit rather than crystallizing into a rank-1 projector.

\section{Gradient Anatomy}
\label{app:gradient}

We analyze gradient norms from training histories across $K \in \{3, \ldots, 36\}$.

\begin{figure}[h]
  \centering
  \includegraphics[width=\textwidth]{outputs/paper_figures/fig_gradient_anatomy.pdf}
  \caption{\textbf{Gradient anatomy.} (a)~$\|\nabla L\|^2$ shows three phases: initial transient, plateau floor, transition spike. (b)~Fisher proxy $\|\nabla L\|^2/L$ drops during plateau. (c)~Gradient waste during plateau scales with $K$. (d)~Higher LR produces smaller gradients but paradoxically longer transitions (anti-Kramers).}
  \label{fig:gradient-full}
\end{figure}

Cumulative gradient norm during the plateau scales as $\tau^{2.14}$ ($R^2 = 0.994$). At $K{=}30$, 19.3\% of all gradient compute is spent in the plateau phase where loss barely changes. This quadratic scaling implies that the computational waste of the plateau grows disproportionately for harder problems.

\section{Representation Geometry}
\label{app:neural-collapse}

We group last-layer residual stream activations by base string $B$ and track within-$B$ variance (how much representations for different $z$ values differ) and between-$B$ variance (how much $B$-group centroids differ)~\citep{rangamani2025neural}.

\begin{figure}[h]
  \centering
  \includegraphics[width=\textwidth]{outputs/paper_figures/fig_neural_collapse.pdf}
  \caption{\textbf{Representation geometry} at $K{=}20$. (a)~Within-$B$ variance rises then slowly decays; between-$B$ variance peaks early and contracts. (b)~Variance ratio increases monotonically. (c)~Total variance peaks post-$\tau$ then decays under weight decay.}
  \label{fig:neural-collapse-full}
\end{figure}

Traditional neural collapse predicts within-class contraction; we observe the opposite---within-$B$ variance \emph{increases} ($1.15\times$) while between-$B$ variance contracts ($0.61\times$). This is because the model must differentiate the $K$ distinct $(B, z_i)$ representations within each $B$ group, shifting capacity from inter-$B$ to intra-$B$ distinctions.

\section{AdamW Timescale}
\label{app:adamw}

We test whether the composite timescale $\tau_\text{adamw} = B/(\eta\lambda D)$ from~\citet{lyu2025power} predicts $\tau_\text{trans}$.

\begin{figure}[h]
  \centering
  \includegraphics[width=\textwidth]{outputs/paper_figures/fig_adamw_timescale.pdf}
  \caption{\textbf{AdamW timescale.} (a)~$\tau_\text{trans}$ vs.\ $\tau_\text{adamw}$ on log-log axes; the relationship is \emph{inverted} ($\tau_\text{trans} \propto \tau_\text{adamw}^{-1.14}$, $R^2 = 0.70$). (b)~Ratio $\tau_\text{trans}/\tau_\text{adamw}$ spans three orders of magnitude across $K$.}
  \label{fig:adamw-full}
\end{figure}

\begin{table}[h]
\centering
\caption{\textbf{Per-sweep regression} of $\log\tau_\text{trans}$ vs.\ $\log\tau_\text{adamw}$ (Transformer only).}
\label{tab:adamw}
\vspace{4pt}
\small
\begin{tabular}{@{}lccc@{}}
\toprule
Sweep & Slope & $R^2$ & $n$ \\
\midrule
$K$-sweep & $-1.70$ & 0.966 & 9 \\
BS-sweep & $-1.16$ & 0.908 & 6 \\
LR-sweep & $-0.53$ & 0.814 & 4 \\
\bottomrule
\end{tabular}
\end{table}

$\tau_\text{adamw}$ anti-correlates with $\tau_\text{trans}$ because task difficulty (which scales with $K$ via dataset size $D = 1000K$) dominates optimizer dynamics. The weight-decay timescale $1/(\eta\lambda) = 100{,}000$ steps is much longer than any observed transition, so weight decay acts as a slow regularizer rather than the primary clock.


\end{document}
